% =====================================================================
%  ch01.tex —— 第 1 章 总导引（PDF 5–9 / 原书 p.5–9）
%  §1.1 为什么是现在？ §1.2 什么是代数拓扑？ §1.3 数学严格性 §1.4 附注
%  原书本章无陈述环境（定义/定理），只有 2 个编号公式。
% =====================================================================
\chapter{总导引}

% [原文 p.5]
\section{为什么是现在？}

尽管我们物理系有不少数学物理学家，研究生院却一直没有数学物理的课程，至少自我大约十年前受聘以来就是如此。我一直想开一门这样的课，但我知道，要在课程表中新增一个带有新科目名称的授课位，需要先跨过许多官僚手续。

去年，我注意到其实已经有一个名为“数理物理学”（mathematical physics）的授课位，不知何故已有约二十年未被使用。结果发现，重新启用一个长期闲置的授课位几乎不需要什么文书工作，于是我决定就这么办。

数学物理可以指很多东西。它常与“物理学家的数学方法”相区别（后者在日语中有不同的译名——物理数学）。它常指理论物理中讨论在数学上严格的那部分，例如统计物理中严格证明热力学极限或相变存在性的那部分。它也可以指理论物理中大量使用各种数学概念（尽管并不十分严格）的那部分，例如弦论。

% [原文 p.6]
我不确定其他教员今后会如何使用这个授课位，但我今年的目标是提供一门面向物理学家的代数拓扑导引，因此它更接近一门“物理学家的数学方法”课程，尽管所覆盖的数学子学科与通常的那些（微积分、线性代数、一点群论等）有所不同。我的理由如下。

正如你们已经知道的，数学在物理中扮演非常重要的角色，Wigner 在其文章 [Wig60] 中对此有著名的论述。但显然并非数学的所有子领域都同等重要。明显重要的有：

\begin{itemize}
  \item 微积分：某种意义上，这门学科源自物理（牛顿等人）。
  \item 线性代数：凡涉及一阶近似的内容都需要它。它也是量子力学（QM）的关键组成部分。
  \item 希尔伯特空间理论：对量子力学同样重要。
  \item 群论：对称性是物理中的基本概念之一。例如，它在分析力学与量子力学中都给出守恒荷。
  \item 微分几何：这对广义相对论（GR）很重要。
\end{itemize}

但也有一些子领域，其用处相当可疑（或至少不是一眼可见）：

\begin{itemize}
  \item 数论：物理主要基于 $\R$ 与 $\C$，而数论研究的是 $\Z$。
  \item 代数几何：它研究由多项式方程定义的对象的形状。物理学家为什么要在意多项式方程？
  \item 数理逻辑/集合论：我们会在物理中用到 Gödel 不完备定理吗？物理会依赖于集合论公理的具体选择吗？
\end{itemize}

直到大约十五年前，代数拓扑都处在这两个清单之间。从 1970 年代起，基础同伦论被用于粒子物理与凝聚态物理中拓扑孤子的研究。从 1980 年代起，一些示性类被用于粒子物理中反常的研究以及凝聚态物理中量子霍尔效应的研究。但也就仅此而已。从数学家的角度看，所用到的代数拓扑也相当初等：物理侧用到的一切大约是 1940 年代发展起来的。因此，代数拓扑通常并不（也不必）包含在典型的物理课程或物理学家的数学方法教科书中。

但过去十五年里情况变了。拓扑绝缘体与超导体成为热门话题，很快我们了解到它们由 $K$ 与 $KO$ 理论分类 [SRFL08, Kit09, RSFL09]。而这是对非相互作用相的分类。对相互作用相（常称为对称保护拓扑（SPT）相或可逆相）的研究，例如 [FK09, CGLW11, GW12, MFCV14]，很快使人意识到分类将由配边群 [KTTW14] 或更一般的上同调理论 [Kit15] 给出。

\takeaway{本节解释这门课为何现在开：拓扑绝缘体、SPT 相与反常用到 $K$、$KO$ 理论、配边群与上同调，而这些内容以往不在物理课程中。课程据此定位为“物理学家的数学方法”之补充。}

% [原文 p.7]
这一预期后来由 [FH16, Yon18] 证实，其中证明 SPT 相的分类由配边群的 Anderson（或 Pontryagin）对偶给出。这也隐含地给出了粒子物理侧关于反常的一般理论。要理解这一切，需要比以往必要程度远为先进的代数拓扑（尽管从数学家的角度看仍不太现代，因为它只用到 1970 年左右的代数拓扑）。

要理解这一迷人的发展，需要我从各种来源学习代数拓扑的零碎知识。但并没有一个地方汇集了面向物理学家的大部分相关材料。这套讲义就是我试图提供这样一个地方的努力。

\section{什么是代数拓扑？它有用吗？}

拓扑是数学中研究空间的子领域。空间当然出现在物理中。我们身处并用广义相对论研究的四维时空是一个数学空间。刚体的位形空间也是一个空间。类似地，凝聚态系统的序参量是一个空间。定义在格点上的所有可能的、有能隙的凝聚态哈密顿量构成的空间也是一个空间。因此，研究空间的方法（即拓扑学）应当对我们有用。

在代数拓扑中，我们通过给空间赋予代数对象来研究空间。给定空间 $X$，与之关联的代数对象有：

\begin{itemize}
  \item \textbf{同伦群}：$\pi_n(X)$ 是通过考虑映射 $S^n\to X$ 得到的群，其中 $S^n$ 是 $n$ 维球面（约定通常的球面是 $S^2$）。
  \item \textbf{配边群}：$\Omega_n(X)$ 是通过考虑映射 $M_n\to X$ 得到的群，其中 $M_n$ 是 $n$ 维的一般光滑流形。
  \item \textbf{同调群}：$H_n(X)$ 是通过考虑 $X$ 中 $n$ 维的“闭链”（本质上是多面体）得到的群，其中闭链是比球面或流形更一般的概念。
\end{itemize}

注意，上述描述中的“源空间”之间有包含关系
\[
  \{\text{球面}\}\subset\{\text{流形}\}\subset\{\text{闭链}\}
  \tag{1.1}
\]
而且不知何故，随着源空间变得更一般，计算的困难程度下降，因此同伦群最难计算，同调群最易计算。

同伦群可用于描述拓扑孤子，配边群在拓扑相以及量子场论反常的研究中起重要作用。我们还会遇到 $K$ 理论 $K(X)$，其计算与定义的困难程度介于同伦与同调之间，与配边群相仿，它在非相互作用拓扑相以及量子场论中费米子场性质的研究中有用。
% [Tutor Note: 原文此处作 “between homotopy and homotopy”，第二个 homotopy 应为 homology；正文按最合理读法译为“同伦与同调之间”，已登记进原文笔误表。]
同调群在此刻的讨论中与物理没有直接关系，但作为最容易计算的对象，得到它们是计算其他代数对象的第一步。

\takeaway{本节给出代数拓扑的基本手法：给空间 $X$ 赋予代数对象再作计算。三个来源空间（球面、流形、闭链）逐级放宽，计算难度则下降，对应同伦群、配边群、同调群；物理上分别用于拓扑孤子、拓扑相与反常。}

% [原文 p.8]
另一个主题是研究空间 $X$ 上的纤维丛 $E$。这是一个带有投影 $E\to X$ 的空间 $E$，使得在每一点 $x\in X$ 附近，$E$ 看起来像 $U\times F$，其中 $x\in U\subset X$ 是 $X$ 的一个邻域，$F$ 是一个固定的空间，称为纤维。纤维丛的拓扑常常借助示性类来区分，这同样是代数拓扑的主题。这一点在物理中出现在许多地方。三个例子：

\begin{itemize}
  \item 规范场（例如描述电磁相互作用或强相互作用的 Maxwell 场或 Yang--Mills 场）在数学上由纤维丛上的联络描述，其中底空间 $X$ 是我们的时空。该纤维丛的拓扑给出 Dirac 磁单极子与瞬子，它们分别由纤维丛的第 1 与第 2 Chern 类这些示性类刻画。
  \item 考虑凝聚态背景下的晶体系统。该系统的能带结构确定布里渊区上的一个（希尔伯特空间的）纤维丛。它的拓扑是该系统许多拓扑性质的根源。例如对 $2+1$ 维系统，该丛的第 1 Chern 类（一个整数）给出量子化的霍尔电导，正如 Thouless、Kohmoto、Nightingale 与 den Nijs 的著名分析所示 [TKNdN82]。
  \item 参数化的量子态族也可以看作参数空间上的一个纤维丛，其纤维可以是波函数的希尔伯特空间，也可以是密度矩阵的空间。该纤维丛的拓扑给出 Berry 相的概念。我们还会看到这样的情况：可以有一族连续的密度矩阵，却永远不能实现为波函数的统计混合的连续族。这一问题同样由一个示性类探测，称为 Dixmier--Douady 类。
\end{itemize}

纤维丛的示性类常取底空间 $X$ 的上同调群中的值，这也解释了（上）同调群在物理中的用处。

本系列讲义的目的是提供最少量的信息，使你们能理解上面几段的内容。不幸的是，这需要整整一个学期。

\section{本课程中的数学严格性}

这不是一门面向（未来的）数学家的课程。对我来说，数学就像一堆有用的应用（在你 Mac/PC 或移动设备上的 app）。下面是一张对照表：

\begin{center}
\begin{tabular}{|c|c|}
\hline
数学 & app \\ \hline
定义 & app 的简短用法说明 \\
定理 & app 本身 \\
证明 & app 的源代码 \\ \hline
\end{tabular}
\end{center}

阅读一个开源 app 的源代码可能有趣且有教益，但若你只想使用这个 app，则并非必需。类似地，若你只想用一次这个 app，那么连读它的简短用法说明（比如在应用商店里）都可能太多。在本课程中，我们不会深入 app 的源代码，但我确实打算告诉你们有哪些 app 可用，并至少给出所涉 app 的简短用法说明。

% [原文 p.9]
我应该说，理论物理的论文与教科书并不像数学中那样包装得整齐。在数学中，你可以跳过证明，仍然用其中的定理来计算我们需要的东西。相反，在物理论文中，我们在假设、推导与结论之间自由地来回穿梭，几乎总是被要求读完整篇论文才能理解其中发生了什么，并将其用于我们自己的目的。不仅如此，数学中的证明通常可靠，所以你可以放心使用定理，不必过分担心。而物理中的陈述常依赖未写明的假设，因此当我们想使用其他论文的结果时需要极其谨慎。也许在这方面，我们物理学家可以从数学家那里学到一些东西……

\takeaway{本课程按数学家的标准组织陈述：定义、定理、证明齐备，但只讲用法、不深究证明——把数学当成 app，重点是知道有哪些工具可用、怎样调用它们。}

\section{附注：“无用”数学分支在物理中的某些罕见用处}

在进入讲义的主要部分之前，我想提一下我在 §1.1 中提到的那些“无用”数学分支与物理之间的一些微薄联系。

\para{数论：}模函数出现在二维共形场论（CFT）与弦论的研究中，并被使用。模函数在数论中同样扮演非常重要的角色。例如，今年出版了一本三枝（Mieda）[三24] 的数论入门书，我常在弦论与二维 CFT 中见到的那些函数，在整本书中处处出现。弦论算不算物理可以争论，但我认为二维 CFT 肯定是物理。所以数论与物理之间至少存在微小的联系。这是偶然吗？还是有更深的关系？我应该说，甚至有一本期刊名为 \textit{Number Theory and Physics}。

\para{代数几何：}在代数几何中，人们研究由多项式方程确定的空间的形状。在物理中，空间的形状通常由微分方程确定。有没有什么地方多项式方程会自然地出现在物理中？如果你好心地把弦论当作物理，那么答案是有的。在弦论中时空是十维的。要描述我们的四维世界，需要假设多出的 $10-4=6$ 维被紧化。真实世界不是超对称的，但如果我们把超对称作为熟悉该系统的途径，那么空间的紧化部分由六维 Calabi--Yau 流形描述。有一个数学定理说，Calabi--Yau 流形总是由多项式方程描述。

\para{数理逻辑/集合论：}在凝聚态系统中，一个标准的玩具模型是自旋链：我们有以整数 $i=-3,-2,-1,0,1,2,3,\dots$ 标记的格点，每个格点上有一个量子比特 $\ket{\pm}_i$。我们考虑由局域相互作用构成的平移不变哈密顿量，例如标准的 Ising 模型：
\[
  H=\sum_i\bigl[a(\sigma_X)_i+b(\sigma_Z)_i(\sigma_Z)_{i+1}\bigr].
  \tag{1.2}
\]

% [原文 p.10]
这类系统以表现出各种有趣现象而著称，研究它已经耗费了无数人时。

一个基本问题是基态是否有能隙。在上面的 Ising 模型中，当 $|a|\neq|b|$ 时有能隙；当 $|a|=|b|$ 时无能隙（并流向共形场论）。可以问：是否存在一种计算机算法，能判定给定的局域、平移不变哈密顿量是否给出有能隙的基态？显然，理论物理学家还没有聪明到能做这件事，但我们能否设想在遥远的未来有朝一日这会成为可能？

答案是否定的 [CPGW15a, CPGW15b, BCLPG18]。要点如下。借助所引文献中解释的一种非常巧妙的构造，可以对任意计算机程序（更确切地说，图灵机）写出一个哈密顿量：在 $L$ 个格点的有限链上，当程序在 $L$ 步内停机时基态有能隙，否则没有能隙。因此，若有一种算法能判定给定哈密顿量在无穷长链上是否给出有能隙的基态，同一算法就能判定给定图灵机是否停机。但计算机科学中一条非常基本的定理说，不存在判定给定图灵机是否停机的算法。所以这是不可能的。

同样，容易写出一个程序，它停机当且仅当当前的标准集合论公理不一致。我们只需枚举所有可能的字母序列：其中许多是垃圾，但有些描述了某个数学定理的证明。任何数学定理的任何可能证明，最终都会出现在枚举过程中。我们让程序在找到 $0=1$ 的证明时停机；而这发生当且仅当当前的标准集合论公理不一致。我们可以把这个程序编码成自旋链哈密顿量。于是这个哈密顿量有有能隙的基态，当且仅当当前的标准集合论公理不一致。而 Gödel 第二不完备定理说，人们既不能证明也不能证伪这个哈密顿量是否给出有能隙的基态。

受这项工作启发，还有若干性质类似的工作。例如 [SM20a, SM20b] 证明，不存在能判断给定量子态是否热化的算法。我还在超对称量子场论（SQFT）中构造了一个例子：不可能存在算法去判定给定的二维 SQFT 是否有超对称真空 [Tac22]\footnote{关于同一意义下不可判定的数学命题，可参看一篇漂亮的科普文章 [Poo12]。}。

这一切在我看来相当迷人。但显然，这并不是数理逻辑／集合论／计算机科学的很深的应用。因为所用到的核心数学事实（Gödel 不完备定理，或图灵机停机问题的不可判定性）虽然根本，却都非常古老。

\takeaway{本章说明为何要为物理学家讲代数拓扑：拓扑绝缘体、SPT 相与反常的分类把同伦群、配边群、$K$ 理论与示性类推到物理前沿。作者交代了课程的定位——讲“有哪些 app 与怎么用”，不深究证明。}

\chapterreview{§1.1 从课程史与物理需求出发，指出代数拓扑在近十五年因拓扑相与反常的分类而变得不可或缺。§1.2 给出三个基本代数不变量：同伦群、配边群、同调群，并用式 (1.1) 的包含关系说明计算难度递减；纤维丛与示性类则是物理中拓扑的另一入口。§1.3 声明课程重在使用而非证明。§1.4 以数论、代数几何、数理逻辑为例，说明所谓“无用”分支与物理的零散联系。全章不涉技术细节，只交代动机与路线。}
